Fine-tuning on narrow tasks can induce \emph{broad} misalignment in a single
language model, and distributed pipelines---federated learning, model merging,
and synthetic-data distillation---create channels by which one contaminated
model can influence many others. We ask whether misalignment spreads through a
population of interacting models the way an infection spreads through a
population of hosts. We answer this in two parts. First, we measure a
misalignment \emph{transmission probability} between real LLMs: exposing a clean
student model to a misaligned teacher's outputs as in-context data, we find a
genuine, dose-dependent contagion in \llama (broad misaligned fraction rising
from $0$ to $16\%$ over $8$ exposures, with full recovery after re-exposure to
aligned data), while \gptmini is effectively immune ($0\%$ at every dose).
Second, we feed the measured transmission ($\beta\!\approx\!0.16$) and recovery
($\mu\!\approx\!0.25$) into a calibrated SIS/SIR-on-networks simulation of $100$
models. The simulation reproduces every classical epidemic signature: a basic
reproduction number \Rzero, a sharp phase transition at the spectral threshold
$R_0=(\beta/\mu)\lambda_{\max}(\mathbf{A})=1$, a topology vulnerability ranking
(scale-free and hub-and-spoke worst), and a critical intervention coverage below
which misalignment becomes endemic. Crucially, the immune-model finding maps onto
herd immunity: a population of ${\sim}70\text{--}80\%$ robust models
self-extinguishes misalignment regardless of topology. We conclude that
misalignment is genuinely contagious and epidemic control is the right frame, but
the highest-leverage lever is per-model robustness, not network structure alone.
